%
%! regex-fix=true
%! dir=chlewey

\songtitle{Hear!}
% -> stories.tex (African Empire -- Monukae); original.tex (Other Narratives)
\tag*{2026}\tag{medieval-ballad}\tag{original-lyrics}\tag{english-lyrics}\tag{african-empire}
\begin{notes}
Written by Carlos Thompson -- a first-person ballad from the \emph{African Empire} alternate history: a Monukae fleet lord who fought the Romans, was knighted by the emperor he served, and came home to see his lineage merge into the Malian imperial line. Retrieved from Carlos Th's Suno page (V6, 3 October 2026).
\end{notes}
\begin{style}
Medieval ballad, close-miked plucked strings with intimate finger and pick detail, measured bass-march pulse, sustained drones beneath kora, lauds and harps, male bass vocals in a resonant, restrained delivery using the user's own recorded voice.
\end{style}
\begin{lyrics}
[Verse 1]
Hear!

I once fought the Romans,
now they knighted me for my service.
They once killed my father—
no blame, that's war on the surface.
Burn in Roman fire
during the last siege in the Aegean Sea,
the same sea my defence
was a duty I fulfilled with glee.

[Verse 2]
Hear!

We defeated the Romans,
but my emperor, my brother-in-arms,
feared my father, his mentor,
and feared me, my fleet and my farms.
I dismantled my rights
to my land and my fleet for his sake;
I put myself to his duty,
and the Roman issue was the highest stake.

[Verse 3]
Hear!

Those Romans who crossed us
were defeated and properly dealt with;
new Romans under Mali
were our new allies and they agreed swith.
I never shifted loyalties—
rather, loyalties shifted around me—
and my beloved Gladiel
rebuilt it while I played diplomacy.

[Verse 4]
Hear!

I took Byzantine Christ
just as I took the Christ of the Franks,
and when Arabs attacked Rome
I gave the Romans my new fleet and my ranks.
They took me as their own;
the emperor knighted Gladiel and me both,
but glee wouldn't endure—
disease found her soon, betraying my troth.

[Verse 5]
Hear!

I left my sons in Byzantium
and came back to Monukae, retired for good.
My first wife also passed,
but my daughter kept the land; she understood.
A lineage of mine goes on,
being merged in the Malian imperial line;
my other lineage keeps the fleet
from the Black Sea to the Caribbean shine.
\end{lyrics}

%! dir=dvigitt

\songtitle{Confusión}
% -> stories.tex (Children's songs and small creatures); original.tex (Other Narratives)
\tag*{2026}\tag{folk-pop}\tag{original-lyrics}\tag{spanish-lyrics}
\begin{notes}
Written by Carlos Thompson -- a barnyard in confusion where each stanza names the same animal by its regional Spanish names: the capybara (\emph{chigüire}, \emph{chigüiro}, \emph{carpincho}), the guinea pig (\emph{curí}, \emph{cuy}, \emph{cobaya}), the turkey (\emph{guajalote}, \emph{pisco}, \emph{guanajo}) and the opossum (\emph{chucha}, \emph{fara}, \emph{tlacuache}). Retrieved from Dvigitt's Suno page (V6-MINI, 3 October 2026).
\end{notes}
\begin{style}
Folk-pop with a vintage 70s mix, lo-fi room ambience, tape saturation, plate reverb, and a low-pass dubstep sub-bass drop; acoustic guitar strum, clapping percussion, accordion stabs, and children’s chorus; bouncing 6/8 waltz; male group vocals trade Spanish narrative lines and playful English call-and-response, building to a festive singalong hook.
\end{style}
\begin{lyrics}
[Verse 1]
El chigüire se levanta
el chigüiro pone la mesa
el carpincho pasa el mate
y el capibara pregunta: what?

[Chorus]
confusión en la granja
confusión en el campo
y todos responden
lo que todos callan

[Verse 2]
un curí pasa corriendo
el cuy dice que no pasa nada
la cobaya mira con recelo
y el conejillo de indias sigue perdido

[Chorus]
confusión en la granja
confusión en el campo
y todos responden
lo que todos callan

[Verse 3]
el guajalote se para altivo
el pisco lo mira de reojo
el guanajo se desentiende
y el pavo los ve y se pierde

[Chorus]
confusión en la granja
confusión en el campo
y todos responden
lo que todos callan

[Verse 4]
la chucha baja del árbol
una fara pasa de largo
el tlacuache se ríe de eso
y la zarigüeya se hace la tonta

[Final Chorus]
confusión en la granja
confusión en el campo
confusión en el monte
confusión en el prado
\end{lyrics}
